%!TEX root = ../main.tex

\section{Conclusion}
\label{section:conclusion}

In this paper, we share our setup for the FACT-AI course at \OurUniversity{}, which teaches FACT-AI topics through reproducibility. 
The course set out to give students 
\begin{enumerate*}[label=(\roman*)]
    \item an understanding of FACT-AI topics, 
    \item an understanding of algorithmic harm, 
    \item familiarity with recent FACT-AI methods, and
    \item an opportunity to reproduce FACT-AI solutions, 
\end{enumerate*}
through a combination of lectures, paper discussion sessions and a reproducibility project. 

Through their projects, our students engaged with the open-source community by creating a public code repository (in the 2019--2020 iteration), as well as with the research community via successful submissions to the MLRC challenge (in the 2020--2021 iteration). 
We also detail how the 2020--2021 iteration brought about its own unique set of challenges due to the COVID-19 pandemic. 

In this course, we illustrate that reproducibility is not only paramount to good science in general, but is also a fundamental component of FACT-AI. 
We received very positive feedback on teaching FACT-AI topics through reproducibility. We believe this was an excellent fit for our students, which not only helped motivate them for the duration of the course, but also helped them develop skills that will be essential in their future research careers, whether in the private or public sector. 
%
%\todo{One sentence on how course generalizes to other topics.}
%
%\MB{Find my suggestion below, it is too long though, we need to make some space for it}
%
To generalize this course setup to other scientific domains, we suggest identifying where the lack of reproducibility in this domain area is coming from and centre the project around evaluating this component. 
%If this is something which could be studied in a 4-week group project based setting (e.g. no large scale user-study), we argue that this course set-up will generalize easily to other domains.